\newcommand{\TabOptimisationParameters}{%
\begin{table}[tb]%
\begin{tabular}{ll}%
	\hline
	Region for optimisation & Approx.\ No.\  of parameters \\
	\hline
	Production target dimensions and location & $3+3$ \\
	Torus1 dipole field strength & 1 \\
	Torus2 dipole field strength & 1 \\
	Muon beam collimator shapes, position, and material & $3+1+1$ \\
	Stopping target shape and location & $4+3$ \\
	Beam blocker position, form, and material& $3+3+1$ \\
	Electron spectrometer dipole field strength& 1 \\
	DIO blockers in the spectrometer & $4$ \\
	\hline
	Approx. total number of parameters& 32 \\
	\hline
\end{tabular}%
	\caption{\tablabel{optimisation:possible-parameters}%
	Aspects of the experiment that can be optimised and estimates for the number of parameters that define each aspect.
	In the case of the target, beam blocker, and collimator shapes the number of parameters is only approximate; crudely speaking there is at least a width, length and height but in principle one could have a very irregular shape that cannot be parametrised by only three numbers, for example, shapes that change as a function of distance along the beamline.
}%
\end{table}%
}

\newcommand{\FigOptimProdTgtLength}{
\begin{figure}[pt]
\centering
\subfloat[][\figlabel{optimisation:ProdTgtSec:Length:Momentum:Muons}Muons]{
\includegraphics[width=0.48\textwidth,trim=0 0 0 1.5cm,clip]{figs/optimisation/ProdTgtGeom/Length_mu-minus_momentum}}
\subfloat[][\figlabel{optimisation:ProdTgtSec:Length:Momentum:Pions}Pions]{
\includegraphics[width=0.48\textwidth,trim=0 0 0 1.5cm,clip]{figs/optimisation/ProdTgtGeom/Length_pi-minus_momentum}}
\caption{\figlabel{optimisation:ProdTgtSec:Length:Momentum}
Change to momentum distributions at the entrance to the first 90 degrees of the bent muon beam solenoid for different target lengths.
Target length is given as half-length which is the Geant4 convention.  
}
\end{figure}
\begin{figure}[pt]
\centering
\subfloat[][\figlabel{optimisation:ProdTgtSec:Length:Integral:Muons}Muons]{
\includegraphics[width=0.48\textwidth,trim=0 0 2cm 2.8cm,clip]{figs/optimisation/ProdTgtGeom/Length_mu-minus_integral_toZero}}
\subfloat[][\figlabel{optimisation:ProdTgtSec:Length:Integral:Pions}Pions]{
\includegraphics[width=0.48\textwidth,trim=0 0 2cm 2.8cm,clip]{figs/optimisation/ProdTgtGeom/Length_pi-minus_integral_toZero}}
\caption{\figlabel{optimisation:ProdTgtSec:Length:Integral}
Integrated muon and pion yields up to a certain momentum at the entrance to the first 90 degrees of the bent muon beam solenoid as a function of target length.
}
\end{figure}
\begin{figure}[pt]
\centering
\subfloat[][\figlabel{optimisation:ProdTgtSec:Length:IntegralRatio:Muons}Muons]{\includegraphics[width=0.48\textwidth,trim=0 0 2cm 2.8cm,clip]{figs/optimisation/ProdTgtGeom/Length_mu-minus_integral_ratios}}
\subfloat[][\figlabel{optimisation:ProdTgtSec:Length:IntegralRatio:Pions}Pions]{\includegraphics[width=0.48\textwidth,trim=0 0 2cm 2.8cm,clip]{figs/optimisation/ProdTgtGeom/Length_pi-minus_integral_ratios}}
\caption{\figlabel{optimisation:ProdTgtSec:Length:IntegralRatio}
Change in the momentum distribution of muons and pions at the entrance to the first 90 degrees of the bent muon beam solenoid as a function of target length.
}
\end{figure}
}

\newcommand{\FigOptimProdTgtRad}{
\begin{figure}[pt]
\centering
\subfloat[][\figlabel{optimisation:ProdTgtSec:Radius:Momentum:Muons}Muons]{\includegraphics[width=0.48\textwidth,trim=0 0 0 1.5cm,clip]{figs/optimisation/ProdTgtGeom/Radius_mu-minus_momentum}}
\subfloat[][\figlabel{optimisation:ProdTgtSec:Radius:Momentum:Pions}Pions]{\includegraphics[width=0.48\textwidth,trim=0 0 0 1.5cm,clip]{figs/optimisation/ProdTgtGeom/Radius_pi-minus_momentum}}
\caption{
\figlabel{optimisation:ProdTgtSec:Radius:Momentum}
Change to momentum distributions at the entrance to the first 90 degrees of the bent muon beam solenoid for different target radii.
}
\end{figure}
\begin{figure}[pt]
\centering
\subfloat[][\figlabel{optimisation:ProdTgtSec:Radius:Integral:Muons}Muons]{\includegraphics[width=0.48\textwidth,trim=0 0 2cm 2.8cm,clip]{figs/optimisation/ProdTgtGeom/Radius_mu-minus_integral_toZero}}
\subfloat[][\figlabel{optimisation:ProdTgtSec:Radius:Integral:Pions}Pions]{\includegraphics[width=0.48\textwidth,trim=0 0 2cm 2.8cm,clip]{figs/optimisation/ProdTgtGeom/Radius_pi-minus_integral_toZero}}
\caption{\figlabel{optimisation:ProdTgtSec:Radius:Integral}
Integrated muon and pion yields up to a certain momentum at the entrance to the first 90 degrees of the bent muon beam solenoid as a function of target radius.
}
\end{figure}
\begin{figure}[pt]
\centering
\subfloat[][\figlabel{optimisation:ProdTgtSec:Radius:IntegralRatio:Muons}Muons]{\includegraphics[width=0.48\textwidth,trim=0 0 2cm 2.8cm,clip]{figs/optimisation/ProdTgtGeom/Radius_mu-minus_integral_ratios}}
\subfloat[][\figlabel{optimisation:ProdTgtSec:Radius:IntegralRatio:Pions}Pions]{\includegraphics[width=0.48\textwidth,trim=0 0 2cm 2.8cm,clip]{figs/optimisation/ProdTgtGeom/Radius_pi-minus_integral_ratios}}
\caption{\figlabel{optimisation:ProdTgtSec:Radius:IntegralRatio}
Change in the momentum distribution of muons and pions at the entrance to the first 90 degrees of the bent muon beam solenoid as a function of target radius.
}
\end{figure}
}

\newcommand{\FigOptimProdTgtFinal}{
\begin{figure}[pt]
\centering
\subfloat[][\figlabel{optimisation:ProdTgtSec:Final:Integral:Muons}Muons]{\includegraphics[width=0.48\textwidth,trim=0 0 0 1.5cm,clip]{figs/optimisation/ProdTgtGeom/OptimalLengthRadius_mu-minus_integral_toZero.pdf}}
\subfloat[][\figlabel{optimisation:ProdTgtSec:Final:Integral:Pions}Pions]{\includegraphics[width=0.48\textwidth,trim=0 0 0 1.5cm,clip]{figs/optimisation/ProdTgtGeom/OptimalLengthRadius_pi-minus_integral_toZero.pdf}}
\caption{\figlabel{optimisation:ProdTgtSec:Final:Integral}
Variation in muon and pion yields as a function of target radius when the total target length is set to the optimised value of 32~cm.
Despite the longer target length the optimal radius is still 1~cm.
}
\end{figure}
}

\newcommand{\FigOptimProdTgtComparePhases}{
\begin{figure}[pt]
\centering
\includegraphics[width=0.9\textwidth,trim=0 0 0 1.5cm,clip]{figs/optimisation/ProdTgtGeom/Plot_compare_phase_1and2.pdf}
\caption{\figlabel{optimisation:ProdTgtSec:Phase1vs2}
Comparison of the muon and pion yields per POT for \phaseI and \phaseII.
The difference arises from the change of target material between the phases.
}
\end{figure}
}
\newcommand{\FigOptimMuBeamDipoleMuStops}{
\begin{figure}[bt]
\centering
%	\fbox{
\includegraphics[width=0.85\textwidth,trim=0 0.5cm 0 1.0cm,clip]{figs/optimisation/MuonBeamDipoles/Tidied_stopped_muons.pdf}
%}
\caption{\figlabel{optim:muBeamDipole:stoppedMu}
	Muon stopping rate as a function of the two dipole field strengths (given relative to the \phaseI design specification).
	A clear anti-correlation is visible which is discussed in the text.
}
\end{figure}
}

\newcommand{\FigOptimMuBeamDipolePiStops}{
\begin{figure}[t]
\centering
%	\fbox{
\includegraphics[width=0.85\textwidth,trim=0 0.5cm 0 1.0cm,clip]{figs/optimisation/MuonBeamDipoles/Tidied_stopped_pions.pdf}
%}
\caption{\figlabel{optim:muBeamDipole:stoppedPi}
	Pion stopping rate as a function of the two dipole field strengths (given relative to the \phaseI design specification).
At the level of statistics used to generate each point, no clear trend is obvious.
Empty squares are those where no pions stopped in the run.
}
\end{figure}
}

\newcommand{\FigOptimMuBeamDipoleMuDispersion}{
\begin{figure}[t]
\centering 
\subfloat[][\figlabel{optim:MuBeamDipole:MuDispersion:Entry}Torus1 Entry]{\includegraphics[width=0.45\textwidth,trim=0 0.9cm 0 1.9cm,clip]{figs/optimisation/MuonBeamCollimators/Tidied_Muon_dispersion_at_entrance.pdf}}
\subfloat[][\figlabel{optim:MuBeamDipole:MuDispersion:Exit}Torus2 Exit]  {\includegraphics[width=0.45\textwidth,trim=0 0.9cm 0 1.9cm,clip]{figs/optimisation/MuonBeamCollimators/Tidied_Muon_dispersion_at_exit.pdf}}
\caption{\figlabel{optim:MuBeamDipole:MuDispersion}
Dispersive effect of the 180\degree bent transport solenoid and dipole field on muons.
No collimating material is yet included, so the high-energy muons being removed is due purely to the beam-pipe itself.
}
\end{figure}
}

\newcommand{\FigOptimMuBeamCollimMuonPaths}{
\begin{figure}[p]
\centering 
	\subfloat[][\figlabel{optim:MuBeamCollim:Beamline:All}All Muons]          {\includegraphics[width=1\textwidth,trim=18cm 1.0cm 26cm 1cm,clip]{figs/optimisation/MuonBeamCollimators/Tidied_All_muons_wGeom.png}}\\
\subfloat[][\figlabel{optim:MuBeamCollim:Beamline:Stopped}Stopped Muons]          {\includegraphics[width=1\textwidth,trim=18cm 1.0cm 26cm 1cm,clip]{figs/optimisation/MuonBeamCollimators/Tidied_Stopped_muons_wGeom.png}}\\
\subfloat[][\figlabel{optim:MuBeamCollim:Beamline:HighP}Muons with $p>70$ MeV/c around the stopping target]          {\includegraphics[width=1\textwidth,trim=18cm 1.0cm 26cm 1cm,clip]{figs/optimisation/MuonBeamCollimators/Tidied_HighP_muons_wGeom.png}}\\
\subfloat[][\figlabel{optim:MuBeamCollim:Beamline:Diff}Stopped $\mu-$High-$p$ $\mu$]{\includegraphics[width=1\textwidth,trim=18cm 1.0cm 26cm 1cm,clip]{figs/optimisation/MuonBeamCollimators/Tidied_Where_to_collimate_wGeom.png}}
\caption{\figlabel{optim:MuBeamCollim:Beamline}
The heights of muons as they pass along the beamline.  
	\protect\subref{fig:optim:MuBeamCollim:Beamline:All} The path of all muons.
	\protect\subref{fig:optim:MuBeamCollim:Beamline:Stopped}: The paths of muons that stop in the target.
	\protect\subref{fig:optim:MuBeamCollim:Beamline:HighP}: The heights of muons with momentum greater than 70 MeV/c when they enter the region around the stopping target.  These could potentially decay in flight to give electrons with 100 MeV/c or greater.
	\protect\subref{fig:optim:MuBeamCollim:Beamline:Diff}: The difference between plot \protect\subref{fig:optim:MuBeamCollim:Beamline:Stopped} and plot \protect\subref{fig:optim:MuBeamCollim:Beamline:HighP}.
	Regions in dark blue would give the greatest impact in removing high-momentum muons whilst leave the stopping muons untouched.
	These plots should be compared to those of \fig{optim:MuBeamCollim:BeamWColl} once collimators have been introduced.
}
\end{figure}
}

\newcommand{\FigOptimMuBeamCollimTransverseSep}{
\begin{figure}[t]
\centering 
\subfloat[][\figlabel{optim:MuBeamCollim:TransverseSep:TS1Entry}At 0\degree]  {\includegraphics[height=0.35\textheight,trim=0.0cm 0.8cm 1.3cm 1.9cm,clip]{figs/optimisation/MuonBeamCollimators/MuonTransversePos_Torus1Tor1FirstPoint}}
\subfloat[][\figlabel{optim:MuBeamCollim:TransverseSep:TS3}At 90\degree] {\includegraphics[height=0.35\textheight,trim=1.7cm 0.8cm 1.3cm 1.9cm,clip]{figs/optimisation/MuonBeamCollimators/MuonTransversePos_Torus2Monitor_0}}
\subfloat[][\figlabel{optim:MuBeamCollim:TransverseSep:TS2Exit}At 180\degree]{\includegraphics[height=0.35\textheight,trim=1.7cm 0.8cm 1.3cm 1.9cm,clip]{figs/optimisation/MuonBeamCollimators/MuonTransversePos_StopTgtSecMonitor_0}}
\caption{\figlabel{optim:MuBeamCollim:TransverseSep}
The separation between stopping and dangerous muons.
The separation is largest at the exit (180\degree), reasonable at the entrance (0\degree), and smallest around the mid-point (90\degree).
}
\end{figure}
}

\newcommand{\FigOptimMuBeamCollimMuonPathsWColl}{
\begin{figure}[ph]
\centering 
\subfloat[][\figlabel{optim:MuBeamCollim:BeamWColl:All}All Muons]                  {\includegraphics[width=1\textwidth,trim=18cm 1.0cm 26cm 1cm,clip]{figs/optimisation/MuonBeamCollimators/Tidied_WColl_AllMuons_WGeom.png}}\\
\subfloat[][\figlabel{optim:MuBeamCollim:BeamWColl:Stopped}Stopped Muons]          {\includegraphics[width=1\textwidth,trim=18cm 1.0cm 26cm 1cm,clip]{figs/optimisation/MuonBeamCollimators/Tidied_WColl_StoppedMuons_WGeom.png}}\\
\subfloat[][\figlabel{optim:MuBeamCollim:BeamWColl:HighP}Muons with $p>70$ MeV/c around the stopping target]          {\includegraphics[width=1\textwidth,trim=18cm 1.0cm 26cm 1cm,clip]{figs/optimisation/MuonBeamCollimators/Tidied_WColl_HighPMuons_WGeom.png}}\\
\caption{\figlabel{optim:MuBeamCollim:BeamWColl}
The heights of muons as they pass along the beamline.  
	\protect\subref{fig:optim:MuBeamCollim:Beamline:All} The path of all muons.
	\protect\subref{fig:optim:MuBeamCollim:Beamline:Stopped}: The paths of muons that stop in the target.
	\protect\subref{fig:optim:MuBeamCollim:Beamline:HighP}: The heights of muons with momentum greater than 70 MeV/c when they enter the region around the stopping target.  These could potentially decay in flight to give electrons with 100 MeV/c or greater.
	These plots should be compared to those of \fig{optim:MuBeamCollim:Beamline} before collimators were introduced, where it is clear how well the dangerous muons are being suppressed.
}
\end{figure}
}

\newcommand{\FigOptimMuBeamCollimTorusOne}{
\begin{figure}[t]
\centering 
\subfloat[][\figlabel{optim:MuBeamCollim:Torus1:perPOT}Particles Surviving per POT]{\includegraphics[width=0.5\textwidth,trim=0.8cm 0.8cm 0.6cm 1.9cm,clip]{figs/optimisation/MuonBeamCollimators/Survived_Coll1_unNormalised-log.pdf}}
\subfloat[][\figlabel{optim:MuBeamCollim:Torus1:fraction}Fraction Surviving Collimator]{\includegraphics[width=0.5\textwidth,trim=0.8cm 0.8cm 0.6cm 1.9cm,clip]{figs/optimisation/MuonBeamCollimators/Survived_Coll1_Normalised-lin.pdf}}
\caption{\figlabel{optim:MuBeamCollim:Torus1}
The effect of changing the height of the collimator in Torus1 on the particle distributions.
}
\end{figure}
}

\newcommand{\FigOptimMuBeamCollimTorusTwoPerPOT}{
\begin{figure}[t]
\centering 
\includegraphics[width=0.95\textwidth,trim=0.3cm 0.3cm 0.8cm 0.1cm,clip]{figs/optimisation/MuonBeamCollimators/Survived_Coll2_unNormalised-lin.pdf}
\caption{\figlabel{optim:MuBeamCollim:Torus2:perPOT}
The number of particles reaching the end of the Torus2 solenoid per POT for different heights of both collimators in Torus1 and Torus2.
}
\end{figure}
}


\newcommand{\FigOptimMuBeamCollimTorusTwoFraction}{
\begin{figure}[t]
\centering 
\includegraphics[width=0.8\textwidth,trim=0.3cm 0.3cm 0.8cm 0.1cm,clip]{figs/optimisation/MuonBeamCollimators/Survived_Coll2_unNormalised-lin.pdf}
\caption{\figlabel{optim:MuBeamCollim:Torus2:fraction}
	The number of particles reaching the end of the Torus2 solenoid relative to the number that enter the Torus1 solenoid (\ie the survival probability) for different heights of both collimators in Torus1 and Torus2.
}
\end{figure}
}

\newcommand{\FigOptimMuBeamCollimTorusTwoContours}{
\begin{figure}[bt]
\centering 
%\fbox{
\includegraphics[width=0.75\textwidth,trim=0.1cm 0.3cm 2.8cm 1.1cm,clip]{figs/optimisation/MuonBeamCollimators/Survived_Coll2_StoppedVsHighP-Muons.pdf}
%}
\caption{\figlabel{optim:MuBeamCollim:Torus2:contours}
Contours showing 2.5 percentage point changes to the stopping (blue) and dangerous (red) muon flux, as a function of the collimator heights.
100\% acceptance is found in the bottom right corner.
%For example, for collimator heights within the first blue contour towards the bottom-right corner, less than 2.5\% of stopped muons are lost.
}
\end{figure}
}

\newcommand{\FigOptimESTDipoleBeamHeightTwoD}{
\begin{figure}[tp]
\centering 
\subfloat[][\figlabel{optim:ESTDipole:Beam:0.0}No Dipole]{
	\includegraphics[width=1\textwidth,trim=5cm 0.5cm 10.3cm 0.3cm,clip]{figs/optimisation/EST_dipole/Tidied_signal_height-dipole_00}}\\
\subfloat[][\figlabel{optim:ESTDipole:Beam:0.1}0.1~T Dipole]{
	\includegraphics[width=1\textwidth,trim=5cm 0.5cm 10.3cm 0.3cm,,clip]{figs/optimisation/EST_dipole/Tidied_signal_height-dipole_10}}\\
\subfloat[][\figlabel{optim:ESTDipole:Beam:0.2}0.2~T Dipole]{
	\includegraphics[width=1\textwidth,trim=5cm 0.5cm 10.3cm 0.3cm,clip]{figs/optimisation/EST_dipole/Tidied_signal_height-dipole_20}}\\
\caption{\figlabel{optim:ESTDipole:Beam}
The heights of electrons along the electron spectrometer that originate in the target with 105~MeV for different dipole field values.
With no dipole field, \protect\subref{fig:optim:ESTDipole:Beam:0.0}, very low energy electrons remain on-axis (straight, blue lines) whilst the signal all drifts vertically and is removed by the beampipe.
At larger dipole fields, such as \protect\subref{fig:optim:ESTDipole:Beam:0.2}, the reverse is true.
}
\end{figure}
}

\newcommand{\FigOptimESTDipoleBeamHeightMean}{
\begin{figure}[tb]
\centering 
%	\fbox{
\includegraphics[width=0.8\textwidth,trim=1.0cm 1.3cm 3.8cm 0.9cm,clip]{figs/optimisation/EST_dipole/Tidied_NoShift-Height}
%}
\caption{\figlabel{optim:ESTDipole:MeanHeight}
Mean height of signal electrons for different values of the dipole field strength.
}
\end{figure}
}

\newcommand{\FigOptimESTDipoleBeamFluxMean}{
\begin{figure}[tb]
\centering 
\includegraphics[width=0.8\textwidth,trim=1.0cm 0.3cm 3.8cm 1.8cm,clip]{figs/optimisation/EST_dipole/Tidied_NoShift-Flux}
\caption{\figlabel{optim:ESTDipole:MeanFlux}
Survival probability for signal electrons as a function of the distance along the beamline for different values of the electron spectrometer's dipole field strengths.
}
\end{figure}
}

\newcommand{\FigOptimESTDipoleAcceptanceVsDipole}{
\begin{figure}[tb]
\centering 
\includegraphics[width=0.6\textwidth,trim=0.3cm 0.3cm 1.9cm 1.1cm,clip]{figs/optimisation/EST_dipole/Tidied_acceptance}
\caption{\figlabel{optim:ESTDipole:acceptance}
Geometric acceptance into the StrECAL detector as a function of the dipole field strength over the electron spectrometer.
}
\end{figure}
}

\newcommand{\FigOptimStopTgtPosMuStops}{
\begin{figure}[bt]
\centering 
%	\fbox{
\includegraphics[width=0.7\textwidth,trim=2.0cm 0.05cm 0.9cm 0.3cm,clip]{figs/optimisation/StopTgtPosition/Tidied_MuonStoppingRate.pdf}
%}
\caption{\figlabel{optim:StopTgtPos:MuStops}
Muon stopping rate per \ac{POT} for different target positions.
The linear behaviour arises from the reduced field strength and fixed target radius such that fewer muons impact the target as it is moved downstream.
}
\end{figure}
}

\newcommand{\FigOptimStopTgtPosSensitivitySpect}{
\begin{figure}[tb]
\centering 
%	\fbox{
\includegraphics[width=0.9\textwidth,,trim=1.3cm 0.2cm 0.7cm 0.6cm,clip]{figs/optimisation/StopTgtPosition/Tidied_-AcceptedMomentum.pdf}
%}
\caption{\figlabel{optim:StopTgtPos:AcceptedMomSpect}
The momentum dependence of the electron acceptance into the detector for different target positions.
The spectrum for each target position is normalised to the muon stopping rate for that position, such that each curve shows the sensitivity to electrons of that momentum.
}
\end{figure}
}

\newcommand{\FigOptimStopTgtPosSensitivityNoBeamBlock}{
\begin{figure}[tb]
\centering 
%	\fbox{
\includegraphics[width=0.9\textwidth,trim=1.3cm 0.2cm 0.7cm 0.6cm,clip]{figs/optimisation/StopTgtPosition/Tidied_NoBeamBlocker-AcceptedMomentum.pdf}
%}
\caption{\figlabel{optim:StopTgtPos:AcceptedMomSpectNoBeamBlock}
The momentum dependence of the electron acceptance into the detector for different target positions when the beam blocker is removed.
}
\end{figure}
}

\newcommand{\FigOptimStopTgtPosSensitivityIntegral}{
\begin{figure}[tb]
\centering 
%	\fbox{
\subfloat[][\figlabel{optim:StopTgtPos:AcceptIntegral:Adjacent}Sensitivity]{%
\includegraphics[width=0.49\textwidth,trim=0.5cm 0.5cm 0.3cm 1.9cm,clip]{figs/optimisation/StopTgtPosition/Tidied_-AcceptedMomentum-Integrated_adjacent.pdf}}
\subfloat[][\figlabel{optim:StopTgtPos:AcceptIntegral:Ratio}Change in Shape]{%
\includegraphics[width=0.49\textwidth,trim=0.5cm 0.5cm 0.3cm 1.9cm,clip]{figs/optimisation/StopTgtPosition/Tidied_-AcceptedMomentum-Integrated_ratios.pdf}}
%}
\caption{\figlabel{optim:StopTgtPos:AcceptIntegral}
\protect\subref{fig:optim:StopTgtPos:AcceptIntegral:Adjacent} 
The variation in sensitivity (acceptance $\times$ stopping rate) to electrons with different momenta as a function of the target position with respect to the nominal location. 
The darkest red line towards the top of the plot represents the sensitivity to signal, and it is that line that should therefore be maximised.
\protect\subref{fig:optim:StopTgtPos:AcceptIntegral:Ratio} 
The change in the shape of the acceptance vs. momentum spectrum as a function of the stopping target location.
}
\end{figure}
}

\newcommand{\FigOptimStopTgtPosHeights}{
\begin{figure}[p]
\centering 
\subfloat[][\figlabel{optim:StopTgtPos:Height:42.5}Electrons from 40 to 45 MeV/c]{%
\includegraphics[width=0.95\textwidth,trim=1.15cm 0.05cm 0.3cm 0.180cm,clip]{figs/optimisation/StopTgtPosition/WithBeamBlocker-Height-VaryShifts-Momentum_42-5.pdf}}%
\\\subfloat[][\figlabel{optim:StopTgtPos:Height:62.5}Electrons from 60 to 65 MeV/c]{%
\includegraphics[width=0.95\textwidth,trim=1.15cm 0.05cm 0.3cm 0.180cm,clip]{figs/optimisation/StopTgtPosition/WithBeamBlocker-Height-VaryShifts-Momentum_62-5.pdf}}%
\\\subfloat[][\figlabel{optim:StopTgtPos:Height:82.5}Electrons from 80 to 85 MeV/c]{%
\includegraphics[width=0.95\textwidth,trim=1.15cm 0.05cm 0.3cm 0.180cm,clip]{figs/optimisation/StopTgtPosition/WithBeamBlocker-Height-VaryShifts-Momentum_82-5.pdf}}%
\\\subfloat[][\figlabel{optim:StopTgtPos:Height:102.5}Electrons from 100 to 105 MeV/c]{%
\includegraphics[width=0.95\textwidth,trim=1.15cm 0.05cm 0.3cm 0.180cm,clip]{figs/optimisation/StopTgtPosition/WithBeamBlocker-Height-VaryShifts-Momentum_102-5.pdf}}%
\\\subfloat[][\figlabel{optim:StopTgtPos:Height:102.5-zoom}Electrons from 100 to 105 MeV/c (Zoomed)]{%
\includegraphics[width=0.95\textwidth,trim=1.15cm 0.05cm 0.3cm 0.180cm,clip]{figs/optimisation/StopTgtPosition/WithBeamBlocker-Height-VaryShifts-Momentum_102-5-Zoom.pdf}}%
\caption{\figlabel{optim:StopTgtPos:Height}
The effect of stopping target position on the height of electrons with a fixed momentum as they pass through the Electron Spectrometer.
The size of the variation indicates the stability of the dipole tune; the two parameters are clearly correlated.
Also striking---particularly in \protect\subref{fig:optim:StopTgtPos:Height:102.5-zoom}---is the way the dependence on the helical pitch angles is affected by the stopping target position.
}
\end{figure}
}

%\newcommand{\FigOptimStopTgtPosSensitivityIntegral}{
%\begin{figure}[tb]
%\centering 
%%	\fbox{
%\includegraphics[width=0.9\textwidth,trim=1.2cm 0.05cm 0.9cm 0.6cm,clip]{figs/optimisation/StopTgtPosition/Tidied_-AcceptedMomentum-Integrated_adjacent.pdf}
%%}
%\caption{\figlabel{optim:StopTgtPos:AcceptedIntegrated}
%The variation in sensitivity to different momentum electrons as a function of the target position, with respect to the nominal location.
%The darkest red line towards the top of the plot represents the signal sensitivity, and it is that line that should therefore be maximised.
%}
%\end{figure}
%}
%
%\newcommand{\FigOptimStopTgtPosSensitivityIntegralRatio}{
%\begin{figure}[tb]
%\centering 
%%	\fbox{
%\includegraphics[width=0.9\textwidth,trim=1.2cm 0.05cm 0.9cm 0.6cm,clip]{figs/optimisation/StopTgtPosition/Tidied_-AcceptedMomentum-Integrated_ratios.pdf}
%%}
%\caption{\figlabel{optim:StopTgtPos:AcceptedIntegratedRatio}
%The relative avvep
%}
%\end{figure}
%}

\newcommand{\FigOptimDIOBeamBlockGeometry}{
\begin{figure}[tb]
\centering 
%\fbox{
\includegraphics[width=0.8\textwidth]{figs/optimisation/BeamAndDIOBlocker/StopTgt_to_Detector-cropped.png}
%}
\caption{\figlabel{optim:DIOBeamBlock:Geometry}
Location of the beam blocker and one possible geometry for the \ac{DIO} blockers, both highlighted in green, shown here before optimisation.
}
\end{figure}
}

\newcommand{\FigOptimDIOBeamBlockESTDispersion}{
\begin{figure}[tb]
\centering 
%\fbox{
\includegraphics[width=0.8\textwidth,trim=0.3cm 0 1.7cm 0.6cm,clip]{figs/optimisation/BeamAndDIOBlocker/Tidied_ElectronDispersion-EST.pdf}
%}\\\fbox{
\includegraphics[width=0.8\textwidth,trim=0.3cm 0 1.7cm 0.6cm,clip]{figs/optimisation/BeamAndDIOBlocker/HandTidied_ElectronDispersion-EST-envelope.pdf}
%\includegraphics[width=0.8\textwidth,trim=1.5cm 0 8.5cm 3.0cm,clip]{figs/optimisation/BeamAndDIOBlocker/Tidied_ElectronDispersion-EST-envelope.png}
%}
\caption{\figlabel{optim:DIOBeamBlock:ESTDispersion}
Momentum-dependent dispersion of electrons passing through the spectrometer.
Top plot: the mean height of different momenta electrons as a function beamline distance, showing how the drift of the centre of gyration is truly proportional to the momentum.
Bottom plot: single standard deviation bands for electrons at different momentum, which shows how the envelope for different momenta overlap considerably, reducing the effectiveness of any collimators.
%Whilst the drift of the centre of gyration is proportional to momentum (top plot), there is reasonable overlap in the overall enevelope
%Mean height of electrons with different momentum from the stopping target to the detector before adding any DIO blockers and before optimised of the beam blocker.
}
\end{figure}
}

\newcommand{\FigOptimDIOBeamBlockAcceptances}{
\begin{figure}[tb]
\centering 
%\fbox{
\subfloat[\figlabel{optim:DIOBeamBlock:Acceptance:40}40 MeV/c]  {\includegraphics[width=0.35\textwidth,trim=0.0cm 0.3cm 0.5cm 1.5cm,clip]{figs/optimisation/BeamAndDIOBlocker/Tidied_Acceptance2D_40.pdf}}\hspace{1cm}
\subfloat[\figlabel{optim:DIOBeamBlock:Acceptance:60}60 MeV/c]  {\includegraphics[width=0.35\textwidth,trim=0.0cm 0.3cm 0.5cm 1.5cm,clip]{figs/optimisation/BeamAndDIOBlocker/Tidied_Acceptance2D_60.pdf}}\\
\subfloat[\figlabel{optim:DIOBeamBlock:Acceptance:80}80 MeV/c]  {\includegraphics[width=0.35\textwidth,trim=0.0cm 0.3cm 0.5cm 1.5cm,clip]{figs/optimisation/BeamAndDIOBlocker/Tidied_Acceptance2D_80.pdf}}\hspace{1cm}
\subfloat[\figlabel{optim:DIOBeamBlock:Acceptance:100}100 MeV/c]{\includegraphics[width=0.35\textwidth,trim=0.0cm 0.3cm 0.5cm 1.5cm,clip]{figs/optimisation/BeamAndDIOBlocker/Tidied_Acceptance2D_100.pdf}}
%}
\caption{\figlabel{optim:DIOBeamBlock:Acceptance}
Acceptance into the straw tracker for electrons with different momentum at the stopping target as a function of the beam and DIO blocker dimensions.
Note the logarthmic scale for the colour bar.
}
\end{figure}
}

\newcommand{\FigOptimDIOBeamBlockHitRate}{
\begin{figure}[tb]
\centering 
%\fbox{
%\hspace{-1.3cm}
\subfloat[\figlabel{optim:DIOBeamBlock:HitRate}Hit Rate from DIO]               {\includegraphics[width=0.45\textwidth,trim=0.8cm 0.3cm 0.3cm 1.05cm,clip]{figs/optimisation/BeamAndDIOBlocker/HitRate2d.pdf}}\hspace{0.04\textwidth}
\subfloat[\figlabel{optim:DIOBeamBlock:HitVAccept}Signal Acceptance vs Hit Rate]{\includegraphics[width=0.45\textwidth,trim=0.8cm 0.3cm 0.3cm 1.05cm,clip]{figs/optimisation/BeamAndDIOBlocker/HitRateVsAcceptance2dRatio.pdf}\vspace{1cm}{}}
%\subfloat[\figlabel{optim:DIOBeamBlock:HitVAccept}Signal Acceptance vs Hit Rate]{\includegraphics[width=0.45\textwidth,trim=0.8cm 0.3cm 2.5cm 1.05cm,clip]{figs/optimisation/BeamAndDIOBlocker/HitRateVsAcceptance2d.pdf}\vspace{1cm}{}}
%\hspace{0.12\textwidth}
%}
\caption{\figlabel{optim:DIOBeamBlock:HitRateAcceptance}
\protect\subref{fig:optim:DIOBeamBlock:HitRate} Number of straw tracker hits per DIO electron.
\protect\subref{fig:optim:DIOBeamBlock:HitVAccept} Ratio between the high-momentum electron ($p>100$~MeV/c) acceptance to the number of hits per DIO electron.  Colour is on a logarithmic scale.
}
\end{figure}
}
